\section{Sistemi Distribuiti con Faults}
\subsection{Faults e Failures}
L'affidabilità totale non esiste nella realtà. Come gestiamo questa cosa?

Prima di tutto, cosa può andare storto?
\begin{itemize}
  \item In base alla causa: fault di esecuzione su un nodo, fault di trasmissione, failure di un componente come un link o un'entità.
  \item In base alla durata: transiente o permanente.
  \item In base all'estensione: localizzato se solo un set (sconosciuto a priori) di componenti del sistema può essere difettoso, 
  diffuso (tutti i nodi)
\end{itemize}

Nessun protocollo è resiliente ad un numero arbitrario di faults.
Nota anche che alcuni fault possono essere più disastrosi di altri, ad esempio i fault non rilevabili rispetto a quelli rilevabili.

Limitiamo il numero e/o il tipo di fallimenti che ammettiamo, per lavorare sotto restrizioni più comode.
Ad esempio, in base a quanto abbiamo speso per quel componente.
L'insieme delle ipotesi sotto le quali lavoriamo si chiama \textbf{failure model}.

Component failure models:
\begin{itemize}
  \item \textbf{link failure}: solo i links possono fallire
  \item \textbf{entity failure}: solo le entità possono fallire
  \item \textbf{hybrid failure}: possono fallire entrambi
\end{itemize}

Considereremo i faults come permanenti: ``once faulty, forever faulty''.

Tipi di entity failures dalla meno alla più pericolosa:
\begin{itemize}
  \item \textbf{crash or fail-stop}: entità che funzionavano correttamente smettono all'improvviso di funzionare
  \item \textbf{send/receive omission}: occasionalmente l'entità perde messaggi in ingresso o messaggi pronti in uscita
  \item \textbf{byzantine fault}: l'entità non segue il protocollo, può eseguire qualsiasi azione, anche \underline{apparentemente} ragionevole
\end{itemize}

Un fallimento bizantino è un qulsiasi fallimento che mostra sintomi diversi a diversi osservatori.

Tipi di byzantine link failures:
\begin{itemize}
  \item \textbf{omission}: il messaggio viene inviato ma mai recapitato
  \item \textbf{addition}: il messaggio è consegnato ad un'entità, anche se nessuno l'aveva generato
  \item \textbf{corruption}: il messagio recapitato è diverso da quello originale
\end{itemize}

I protocolli fault tolerant (FT) dipendono anche dalla topologia del grafo su cui sono applicati.
Problemi non triviali sono risolti con la partecipazione di tutti i nodi. Serve connettività.

\defbox{ \textbf{Edge connectivity}: numero minimo di archi che, se rimossi, disconnettono il grafo }
\defbox{ \textbf{Node connectivity}: numero minimo di nodi che, se rimossi, disconnetono il grafo }

Queste due quantità stanno nella seguente relazione: node connectivity $\leq$ edge connectivity $\leq$ max degree

% Eg. TODO tabella
% network - max deg - node conn - edge conn
% tree    - n-1 (star) - 1 - 1
% bidirectional ring - 2 - 2 - 2
% grafo completo - n-1 - n-1 - n-1

\theobox{ Lemma: Se $k$ \textit{link} arbitrari possono fallire, è impossibile effettuare il broadcast su reti con una edge connectivity minore di $k+1$ }
\theobox{ Lemma: Se $k$ \textit{nodi} arbitrari possono fallire, è impossibile effettuare il broadcast su reti con una node connectivity minore di $k+1$ }

\subsection{Problema del p-Agreement}
Input: ogni entità ha in input un valore $v_x$ da un set conosciuto di valori $S$.\\
Output: almeno $p$ entità devono accordarsi sullo stesso valore $d \in S$ in un tempo finito.

Operiamo in una condizione di non banalità: inizializzati tutti i nodi con lo stesso valore, devono accordarsi su tale valore.

Se $p = n$, con $n$ totalità delle entità, siamo nel caso del \textbf{consenso}.
Ovvero, tutti i nodi si sono accordati su un valore comune.

Il problema del consenso si manifesta in casi reali come nei database distribuiti per decidere se una transazione è da committare o scartare,
diagnostica per decidere se un nodo è diffettoso o funzionante, o voting.

Analizziamo ora il problema del consenso nel caso di reti con \textbf{link failures} di tipo \textit{omission}.
Vincoli di output:
\begin{itemize}
  \item \textbf{argreement}: tutte le entità devono accordarsi sullo stesso valore
  \item \textbf{non banalità}: se tutte venongo inizializzate con lo stesso valore, quello è il valore su cui deve raggiungere
  il consenso
  \item \textbf{terminazione}: in un tempo finito
\end{itemize}

Il racconto denominato \textit{Two general's tale} spiega tale problema (consenso su reti con omission link failures) in modo intuitivo.

\subsection{Two general's tale}
Due generali dello stesso esercito si trovano in cima ad due montagne diverse e devono decidere insieme SE e IN QUALE MOMENTO attaccare, 
per farlo devono scambiarsi dei messaggi. Non possono farlo visivamente perchè verrebbero intercettati.
Si scambiano messaggi attraverso i messaggeri, che passano per una vallata; nella vallata ci sono però gli scout nemici, che potrebbero
uccidere i messaggeri. Una volta inviato un messaggio deve essere restituito un ack. Il problema?\\
Il link (la vallata) è inaffidabile, anche gli ack potrebbero non arrivare. Servirebbe l'ack dell'ack dell'ack ecc.

\theobox{ Questo problema non ha soluzione, neanche in sistemi sincroni. }

\theobox{ \textbf{Lemma}. In qualsiasi esecuzione di un protocollo funzionante in cui si decide di attaccare, almeno un messaggio
deve essere consegnato. Altrimenti, sarebbe impossibile distinguere tra lo scenario in cui un generale ha deciso di non attaccare
e lo scenario in cui ha deciso di attaccare ma il messaggio è stato perso nella vallata. }

Dim per assurdo.\\
Assumiamo che esista un protocollo che risolva il problema e che il link trasporti almeno un messaggio prima del fallimento. \\
Prendiamo un'esecuzione $E$ non banale in cui il protocollo porta entrambi i generali ad attaccare, con il minor numero di messaggi possibile $k \geq 1$. \\
Prendiamo una seconda esecuzione $E'$ in cui l'ultimo messaggio, da A a B, non viene consegnato. \\
Le due istanze sono indistinguibili per A. Dall'assunzione che il protocollo sia corretto sappiamo che
anche B deve aver attaccato (le loro mosse devono essere concordi). \\
Conseguenza: i due generali hanno attaccato con $k-1 \geq 0$ messaggi.
Dividiamo i due casi:
\begin{itemize}
 \item se $k-1 > 0$ abbiamo una contraddizione. La prima esecuione, $E$, doveva essere quella che richiedeva il minor numero di messaggi,
 ma ne abbiamo trovata una con un numero minore (k-1), ovvero $E'$.
 \item se $k-1 = 0$ abbiamo una contraddizione. La seconda esecuzione è andata a buon fine con zero messaggi, ma almeno uno era richiesto.
\end{itemize}

Nota che ciò che fa fallire la dimostrazione è il fatto che il link potrebbe fallire, non il fatto che fallisca o meno effettivamente.

Pensa ad esempio alle spunte blu di Whatsapp in una conversazione tra due persone. Le spunte blu servono per aggirare questo problema e
sono fonite dal server, un terzo attore della comunicazione.

\subsection{Algoirtmi Precedenti con Link failures}
\subsubsection{Consenso con link failures}
\theobox{ Se il numero di fallimenti $F$ supera o paraeggia la edge connectivity, è impossibile raggiungere il consenso con un
protocollo deterministico. La edge connectivity deve essere almeno $F+1$. }

Assumiamo che la edge connectivity sia almeno $F+1$.
Il broadcast continua a funzionare siccome la rete non è completamente disconnessa.

Per esempio, il Flooding continua a funzionare
\begin{itemize}
  \item $\text{Messages}(\text{Flooding}) \leq 2m$
  \item $\text{Time}(\text{Flooding}) \leq D(G)$, con $D(g)$ diametro del grafo con faulty links
\end{itemize}

Idea: gni entità fa broadcast del suo valore, ogni entità applica la stessa funzione sui valori ricevuti, alla fine tutte convergono allo
stesso risultato. Il consenso è possibile.

\subsubsection{Two Steps}
Broadcast in un \underline{grafo completo} con \underline{fallimenti dei link: algoritmo \textbf{Two Steps}}.

Assumiamo ci siano al massimo $F$ fallimenti $< n-1$ e $F$ sia conosciuto al nodo $x$ che inizia il broadcast.\\
Tip. Tecnica dell'avversario. Tecnica che consiste nel ``giocare'' ad essere un avversario del protocollo, utile per analizzare il caso peggiore.

Bastano due passaggi per effettuare il broadcast:
\begin{enumerate}
  \item $x$ manda $I$ a $F+1$ vicini. Sicuramente almeno ad un vicino arriva $I$.
  \item ogni entità ricevente il messaggio da $x$ lo propaga a tutti i vicini, escluso $x$.
\end{enumerate}

Message Complexity: al massimo $(F+1) + (F+1) (n-2) = (F+1) (n-1) \in O(Fn)$ messaggi.\\
La prima parte sono i messaggi mandati da x. La seconda parte sono i messaggi mandati dagli altri nodi. Ogni nodo
del secondo step deve mandare il messaggio a $n-2$ nodi: evita se stesso e quello da cui ha ricevuto $I$

\subsubsection{FT-BroadElection}
Leader Election in Grafo Completo con Link Failures: \textbf{FT-BroadElection}.

Restrizioni: start simultaneo, ID univoci, $F < n-1$ known, conoscenza della topologia (grafo completo).

\textbf{FT-BroadElection}:
\begin{enumerate}
  \item ogni entità manda il suo ID con Two Steps
  \item ogni entità calcola il minimo e decide se è il leader
\end{enumerate}

Costo di Two Steps ripetuto per $n$ volte: $n (F+1) (n-1) \in O(F \cdot n^2)$

\subsection{Entity Faults}
Gli stessi vincoli di prima sono richiesti: agreement per le entità non faulty, non trivialità sulle entità non faulty, terminazione
sulle entità non faulty.

Lavoriamo nel caso in cui possano avvenire sono entity faults.

\theobox{ È impossibile raggiungere il consenso in reti con entità soggette a crash, anche se ne fallisce solamente una, appartenente a
a un grafo totalmente connesso (la migliore situazione possibile), usando protocolli deterministici su sistemi asincroni. }

L'idea della dimostrazione è che i tempi di attesa sulla rete sono finiti ma \textit{impredicibili}, quindi non possiamo distinguere
se una entità è crashata o se deve ancora inviare un messaggio.

Soluzione: aggiungiamo la restrizione di lavorare su \textbf{sistemi totalmente sincroni}, con entity failures di tipo \textbf{crash}.

In questa condizione riusciamo a progettare algoritmi che ammettono fino a $f < n$ crash.

Assunzione: i valori ammissibili su cui chierchiamo il consenso sono $v(x) \in S = \{0, 1\}$

Idea: questi algoritmi lavorano in \textbf{step}. Ad ogni time step, ogni entità non faulty invia un messaggio chiamato
\textbf{report}. Diventa quindi facile distinguere se un'entità è crashata o meno: se non invia nessun messaggio
durante il timestep designato, significa che è avvenuto un crash. Ad ogni timestep:
\begin{itemize}
  \item $x$ invia un report ad ogni $y \in N(x)$
  \item $x$ si aspetta un report da ogni $y \in N(x)$
\end{itemize}

$x$ riesce a rilevare se un suo vicino è crashato al tempo $t$ dal fatto che non riceve un messaggio entro $t+1$.

Indichiamo con $r(x, t)$ il report di $x$ al tempo $t$.\\
Il report è inizializzato a $r(x, 0) = v(x)$

Cosa contiene un report?\\
Nei casi che stiamo prendendo in analisi, solamente uno 0 o un 1.

% Algoritmi di consenso randomizzato.

% Oss. al round r i messaggi myvalue contengono o 0 o 1 e, al massimo un valore tra 0 e 1, può essere quello con la maggioranza assoluta
% (almeno n/2+1).

% Dato che mando un messaggio con un valore specifico, ovvero non "?", solamente se vedo la metà più uno, allora durante un round girano
% solo proposte con lo stesso valore.

% L'algoritmo garantisce non banalità. TODO

% L'algoritmo garantisce agreement. TODO

% L'algoritmo garantisce la terminazione in un numero di round $\in O(2^n)$.
% Abbiamo due scenari possibili:
%  - se un nodo x riceve la maggioranza dei valori v, propone v a tutti gli altri.
%  - se nessuna entità vede la maggioranza allo stadio di proposal, non è detto che termini al round successivo.

% Se al prox round almeno n/2+1 nodi generano lo stesso valore, viene generata una maggioranza.
% La probabilità che quel numero di nodi si accordi sullo stessso valore, è sicuramente maggiore della probabilità che TUTTE le entità
% scelgano lo stesso valore: P_{majority same value} >= P_{all same value} = 1/2^n

% Quindi... se sono termino al round r vuol dire che ho fallito nei r-1 round precedenti e ho avuto successo in quello attuale: p (1-p)^{r-1}.

% Fault Bizantini.
% I fault bizantini sono più difficili da affrontare perchè le enità si comportano in modo diverso nel tempo: possono seguire il protocollo per un pò,
% poi simulare un crash, poi riprendere, ecc.

% Anche gli attacchi informatici in cui l'attaccante prende il controllo di un nodo sono un esempio di comportamenti bizantini.
% Anche la collaborazione tra nodi bizantini è possibile e potrebbe fuorviare il protocollo.

% Teorema. Esiste un risultato secondo il quale è possibile progettare un protocollo deterministico che raggiunga il consenso tollerando
% al più f < n/3 fault bizantini.

% Consenso su sistema sincrono con failure bizantini.
% Restrizioni: sincronicità, grafo completo, connettività completa, failure bizantina, inizio sync, ID univoci, ogni nodo conosce gli IDs
% dei suoi vicini (restrizione aggiuntiva per fare fronte all'impredicibilità dei fault bizantini).

% Mandiamo solo gli zero, gli 1 sono dedotti dalla mancata comunicazione, possibile siccome il protocollo è sincrono.

% Aggiungiamo ai messaggi l'ID del proposer e il timestamp di invio del messaggio: (0, ID(x), t).
% In questo modo possiamo rilevare un forge (riceviamo msg da ID doppi o da nodi che non sono nostri vicini), e un tempo fittizio (sistema sincrono).

% Oss. L'importante è che anche i nodi bizantini si comportino in modi coerente, rispettando il protocollo. Rileviamo solo comportamenti
% "bizzarri"; e.g. far accettare ad alcuni lo zero, ad altri un uno. Non possiamo impedire che il bizantino proponga un valore iniziale
% a suo piacere.

% Vediamo il meccanismo RegisteredMail per comunicare lo zero. Assomiglia al broadcast con faulty link.
% RegisteredMail per nodo non faulty x:
% Posso effettuare un wakeup mandando (init, 0, ID(x), t). \\
% Ricevuto un messaggio (init, 0, ID(y), t) da un nodo y facciamo:
%  - broadcast di (echo, 0, ID(y), t) se il messaggio proviene dal canale su cui è registrato y, al tempo t+1, ed è il primo init ricevuto da quell'ID.
%  - ignoriamo il messaggio in tutti gli altri casi. 3 scenari: tempo falso, ID falso wrt link registrato, init doppio

% Dal teorema, almeno f+1 non faulty echos arrivano. TODO: completa ragionamento.

% Se un bizantino forgia un messaggio, potrebbe accordarsi con gli altri bizantini per propagare questo messaggio senza rispettare il protocollo.
% I nodi non faulty vengono ingannati? No, non dovrebbero. Anche assumendo che tutti i bizantini facciano eco del messaggio,
% i bizantini sono uno in meno rispetto ai non faulty: f < f+1.

% TODO: 1a ora lezione 20 novembre
% Prima cosa spiegata: TellZero_byzantine

% Message Complexity of TellZero_Byz.
% n-1. Un messaggio a tutti
% Extra per esame: consenso randomizzato con fault bizantini, su sistemi asincroni.
